\chapter*{Conclusión}
\addcontentsline{toc}{chapter}{Conclusión}

El Bloque I puede condensarse en una sola cadena conceptual:

\begin{center}
\begin{tcolorbox}[
  width=0.48\linewidth,
  colback=grisclaro,
  colframe=gris,
  boxrule=0.6pt,
  arc=1.5mm
]
\centering
estructura real\\
$\downarrow$\\
modelo\\
$\downarrow$\\
cargas y vínculos\\
$\downarrow$\\
equilibrio\\
$\downarrow$\\
$N,\ V,\ M,\ T$\\
$\downarrow$\\
$\sigma,\ \tau$\\
$\downarrow$\\
$\varepsilon$\\
$\downarrow$\\
$\vect u$
\end{tcolorbox}
\end{center}

Dominar esta cadena es mucho más importante que memorizar fórmulas aisladas.

Cuando llegue el resto de la asignatura ---axil, flexión, torsión, cortante, energía y elasticidad tridimensional--- no aparecerá una teoría completamente nueva: se irán añadiendo herramientas a esta misma estructura lógica.
